\documentclass[11pt,a4paper]{article}
\usepackage[margin=25mm]{geometry}
\usepackage{amsmath,amssymb,bm,newtxtext,newtxmath,booktabs,microtype}
\usepackage[hidelinks]{hyperref}
\newcommand{\Tr}{^{\mathsf T}}
\newcommand{\ip}[2]{\langle #1,#2\rangle_{\mathrm F}}
\title{A Consistent Forward--Adjoint Orbital Metric}
\author{Implementation note}
\date{16 October 2026}
\begin{document}
\maketitle

\section{Physical tangent}
Let $A$ and $B$ contain the inactive and active raw orbital coefficients in
the atomic-orbital (AO) basis, and let $S$ be the AO overlap matrix.
All quantities without a variation symbol are fixed at the accepted point.
Define
\begin{align}
 H_i&=A\Tr SA=R_i\Tr R_i, & S&=L\Tr L,\\
 P&=I-AH_i^{-1}A\Tr S, & D&=PB, & K&=H_i^{-1}A\Tr SB.
\end{align}
Here $R_i$ and $L$ are upper-triangular Cholesky factors, $P$ is the inactive
complement, and $D$ contains the projected active orbitals. Given allowed
raw-coefficient variations $(\delta A,\delta B)$, the physical derivatives are
\begin{align}
 h&=P\delta A,\\
 v&=P\delta B-hK-AH_i^{-1}h\Tr SB. \label{eq:v}
\end{align}
The last term in Eq.~\eqref{eq:v} is part of the changing inactive projector;
it is not an optional structure response.

For an individual active ray $d_j=D[:,j]$, write
\begin{equation}
 n_j=d_j\Tr Sd_j,\qquad Q_j=I-d_jd_j\Tr S/n_j.
\end{equation}
The tangent feature map $\mathcal J$ is
\begin{equation}
 F_i=LhR_i^{-1},\qquad f_{a,j}=LQ_jv_j/\sqrt{n_j},\qquad
 \mathcal J(\delta A,\delta B)=(F_i,F_a).
 \label{eq:feature}
\end{equation}
It measures inactive-subspace motion and normalized active-ray motion.
This ray representation applies to selected valence-bond structure spaces;
it assumes no closure under active-orbital transformations.

Only when the orbital manifold has the complete full-AO active-space gauge
may the active rays be replaced by their span. In that case define
\begin{equation}
 H_a=D\Tr SD=R_a\Tr R_a,\qquad Q_a=I-DH_a^{-1}D\Tr S,
 \qquad F_a=LQ_avR_a^{-1}. \label{eq:span}
\end{equation}
Completing a computational spin-string carrier does not establish this
physical gauge for an arbitrarily selected structure space.

\section{Strict adjoint}
Let $Z_i$ and $Z_a$ be feature covectors. Reverse precisely the operations in
Eqs.~\eqref{eq:feature}--\eqref{eq:span}:
\begin{align}
 W_i&=L\Tr Z_iR_i^{-\mathsf T},\\
 w_{a,j}&=Q_j\Tr L\Tr z_{a,j}/\sqrt{n_j}
 \quad\text{(individual rays)},\\
 W_a&=Q_a\Tr L\Tr Z_aR_a^{-\mathsf T}
 \quad\text{(complete active span)}.
\end{align}
The raw-coefficient covectors are then
\begin{align}
 \overline B&=P\Tr W_a,\\
 \overline A&=P\Tr\left[W_i-W_aK\Tr-SBW_a\Tr(AH_i^{-1})\right].
 \label{eq:adjoint}
\end{align}
These equations follow by the Frobenius identity
\begin{equation}
 \ip{Z_i}{F_i}+\ip{Z_a}{F_a}
 =\ip{\overline A}{\delta A}+\ip{\overline B}{\delta B}.
\end{equation}
In particular, the last term in Eq.~\eqref{eq:adjoint} uses the transpose of
the \emph{stored forward factor} $AH_i^{-1}$. It must not independently
reconstruct that transpose by assuming an explicitly computed inverse is
exactly symmetric in floating-point arithmetic.

\section{Metric action and ownership}
If $U$ denotes the existing nonredundant sparse lift, the reduced metric is
\begin{equation}
 G=U\Tr\mathcal J^*\mathcal JU,\qquad
 p\Tr Gp=\|\mathcal JU p\|_{\mathrm F}^2. \label{eq:metric}
\end{equation}
Apply the forward feature map and its strict adjoint using the same cached
accepted-point factors. The norm and action then describe one computational
graph, instead of separate expressions related by projector identities such
as $P\Tr SP=SP$ that need not hold to working precision for computed factors.
No reduced metric matrix, inverse, eigendecomposition, or coordinate
whitening is introduced by Eq.~\eqref{eq:metric}. Its application has shared
$O(N_{\mathrm{BF}}(n_i+n_a))$ tangent/cotangent storage and leading work
$O(N_{\mathrm{BF}}^2(n_i+n_a)+N_{\mathrm{BF}}(n_i+n_a)^2)$.
The existing accepted-point AO and small orbital factors are reused; no
large array is replicated per OpenMP worker.
Matrix products are explicitly grouped to apply each AO operator to the
narrow tangent or cotangent first: the implementation never forms $LQ_a$ or
$Q_a\Tr L\Tr$ during an action. Such products would introduce avoidable
$O(N_{\mathrm{BF}}^3)$ work despite the factorized mathematical formula.

Literal bitwise symmetry is not asserted. This change removes a distinct
forward/adjoint graph mismatch, without changing the AH symmetry certificate,
orbital coordinates, Hessian--vector product, or convergence tolerances.
The independent diagnostic constructs the feature Gram matrix at a captured
molecular point; that dense matrix is a diagnostic only, never a production
optimizer representation.

\begin{center}
\begin{tabular}{ll}
\toprule
Operation & Implementation \\
\midrule
Forward map & \texttt{physical\_feature()} \\
Strict adjoint & \texttt{feature\_adjoint()} \\
Metric composition & \texttt{apply()} \\
\bottomrule
\end{tabular}
\end{center}
The previous expanded action and its unused caches are deleted.
\end{document}
